\documentclass[11pt,a4paper]{article}
\usepackage[T1]{fontenc}
\usepackage[utf8]{inputenc}
\usepackage{lmodern,microtype,amsmath,amssymb,mathtools,amsthm,booktabs}
\usepackage[margin=25mm,headheight=16pt]{geometry}
\usepackage{hyperref,xcolor,fancyhdr,graphicx}
\definecolor{accent}{HTML}{9B4635}
\hypersetup{colorlinks=true,linkcolor=accent,urlcolor=accent,citecolor=accent,pdftitle={Formalism: Griffiths Chapter 3 Solutions}}
\numberwithin{equation}{section}
\pagestyle{fancy}\fancyhf{}\lhead{\small Quantum Mechanics / Chapter 3: Formalism}
\rhead{\small Kumar Azad}\cfoot{\thepage}
\newcommand{\dd}{\,\mathrm{d}}
\newcommand{\ii}{\mathrm{i}}
\newcommand{\ee}{\mathrm{e}}
\newcommand{\ket}[1]{\left|#1\right\rangle}
\newcommand{\bra}[1]{\left\langle#1\right|}
\newcommand{\braket}[2]{\left\langle#1\,\middle|\,#2\right\rangle}
\newcommand{\avg}[1]{\left\langle#1\right\rangle}
\newcommand{\comm}[2]{\left[#1,#2\right]}
\newcommand{\problem}[2]{\subsection*{Problem #1 — #2}\addcontentsline{toc}{subsection}{Problem #1 — #2}}
\newcommand{\qnote}[1]{\par\smallskip\noindent\textit{Verification / interpretation.} #1\par\smallskip}
\newcommand{\answer}[1]{\par\smallskip\noindent\textcolor{accent}{\textbf{Result:}}\quad #1\par\smallskip}
\setlength{\parindent}{0pt}\setlength{\parskip}{5pt}
\title{\textcolor{accent}{\small MATHEMATICAL PHYSICS · EDUCATIONAL COMPANION}\\[8pt]
\textbf{\Huge Quantum Mechanics:\\Formalism}\\[11pt]
{\large Independent worked solutions to Chapter 3}\\[4pt]
{\normalsize Problems 3.1–3.49}}
\author{Kumar Azad\\ \small\url{https://krazad0.github.io/QM-CH1-SOLUTION/chapter3/}}
\date{November 2026}
\begin{document}
\maketitle
\begin{abstract}
Independent, original stepwise calculations for the forty-nine numbered exercises
in Chapter 3 (\emph{Formalism}) of Griffiths and Schroeter,
\emph{Introduction to Quantum Mechanics}, third edition.
The material develops Hilbert spaces, Hermitian operators, continuous
spectra, Fourier-dual representations, uncertainty relations, Dirac
notation, matrix mechanics, coherent states, supersymmetric
Hamiltonians, and operator domains. Short descriptions are
paraphrased, not copied from the textbook; consult the original for
the full questions. This is an academic-style coursework
companion, \textbf{not an arXiv submission or a peer-reviewed paper}.
\end{abstract}
\tableofcontents
\section{Notation, domains and inner products}
For position-space wavefunctions, the Hilbert space is
$L^2(\mathbb R)$ or $L^2(a,b)$, as appropriate, and
$\braket{f}{g}=\int f^*(x)g(x)\dd x$.
A physical observable is represented by a self-adjoint operator
on a specified domain; checking formal integration by parts
is generally only a symmetry check. For sufficiently regular
states and a real time-independent potential,
\begin{align}
\hat H&=\frac{\hat p^2}{2m}+V(\hat x),&
\hat p&=-\ii\hbar\frac{\dd}{\dd x},&
\comm{\hat x}{\hat p}&=\ii\hbar I,\\
\frac{\dd}{\dd t}\avg Q&=
\frac{\ii}{\hbar}\avg{\comm{H}{Q}}+
\left\langle\frac{\partial Q}{\partial t}\right\rangle,&
\sigma_A^2&=\avg{A^2}-\avg A^2.
\end{align}
The normalized momentum wavefunction is
$\Phi(p)=(2\pi\hbar)^{-1/2}\int\ee^{-\ii px/\hbar}\Psi(x)\dd x$.
Where eigenvectors are continuous-spectrum generalized states,
Dirac deltas replace Kronecker deltas.
\section{Hilbert space and Hermitian operators (3.1–3.10)}
\problem{3.1}{Square-integrable functions form a vector space and inner-product space}
Let $L^2(a,b)$ denote equivalence classes of measurable complex functions
with finite $\|f\|^2=\int_a^b|f|^2\dd x$.
For $f,g\in L^2$ and complex $\alpha,\beta$, use
$|u+v|^2\le2|u|^2+2|v|^2$ to obtain
\begin{equation}
\|\alpha f+\beta g\|^2\le
2|\alpha|^2\|f\|^2+2|\beta|^2\|g\|^2<\infty.
\end{equation}
Thus the set is closed under vector addition and scalar multiplication;
the other vector-space axioms hold pointwise almost everywhere.
The \emph{normalized} subset is not a vector space: it does
not contain zero and, for normalized $f$, the multiple $2f$
has norm two.
For $\braket f g=\int f^*g$, linearity in the second slot gives
$\braket f{\alpha g+\beta h}=\alpha\braket fg+
\beta\braket fh$. Conjugate symmetry follows by complex
conjugation under the integral; positivity is
$\braket ff=\int|f|^2\ge0$, with equality exactly
for the zero equivalence class. The Cauchy--Schwarz
inequality assures finite overlaps.
\answer{$L^2(a,b)$ is an inner-product vector space and is complete
as a normed space; the unit sphere of normalized states is not
a linear subspace.}

\problem{3.2}{Integrability and the unbounded derivative operator}
For $f(x)=x^\nu$, $0<x<1$, square integrability requires
\begin{equation}
\int_0^1|f|^2\dd x
=\int_0^1x^{2\nu}\dd x
=\frac{1}{2\nu+1}<\infty
\quad\Longleftrightarrow\quad\boxed{\nu>-\frac12}.
\end{equation}
For $\nu=1/2$, $f=\sqrt x$ and $xf=x^{3/2}$
are in $L^2(0,1)$, with squared norms $1/2$ and $1/4$.
But $f'=1/(2\sqrt x)$ has
$\int_0^1|f'|^2\dd x=\frac14\int_0^1\dd x/x=\infty$.
\answer{Differentiation does not map every $L^2$ function
back into $L^2$; the derivative operator requires a domain.}

\problem{3.3}{Equivalence of Hermiticity tests via polarization}
Suppose $\braket h{Qh}=\braket{Qh}h$ for all
$h$ in a complex linear domain.
Set $B(f,g)=\braket f{Qg}-\braket{Qf}g$.
Then $B(h,h)=0$, while $B$ is conjugate-linear
in its first entry and linear in its second.
Expanding $B(f+g,f+g)=0$ gives
\begin{equation}
B(f,g)+B(g,f)=0.
\end{equation}
Expanding $B(f+\ii g,f+\ii g)=0$ instead gives
$\ii B(f,g)-\ii B(g,f)=0$. Together they imply
$B(f,g)=0$ for arbitrary $f,g$:
\begin{equation}
\boxed{\braket f{Qg}=\braket{Qf}g}.
\end{equation}
The converse follows by setting $f=g$.
\qnote{This proves the algebraic symmetry identity;
self-adjointness additionally requires equality of adjoint domains.}

\problem{3.4}{Sums, scalar multiples and products of Hermitian operators}
On an appropriate common dense domain,
$(A+B)^\dagger=A^\dagger+B^\dagger$;
thus the sum of symmetric operators is symmetric.
For a scalar $\alpha$ and nonzero Hermitian $Q$,
$(\alpha Q)^\dagger=\alpha^*Q$, so
$\alpha Q$ is Hermitian precisely when $\alpha$ is real.
For Hermitian $A,B$, $(AB)^\dagger=BA$;
$AB$ is Hermitian when $AB=BA$ (with the
necessary shared operator-domain conditions).
Position is multiplication by real $x$:
$\int f^*xg=\int(xf)^*g$.
For $H=-\hbar^2\partial_x^2/(2m)+V(x)$,
two integrations by parts give
\begin{equation}
\braket f{Hg}-\braket{Hf}g
=-\frac{\hbar^2}{2m}
\left[f^*g'-f'^*g\right]_{\partial I},
\end{equation}
where $V$ is real. The surface term vanishes
for suitable boundary conditions, proving formal symmetry.
\qnote{Sums and products of unbounded self-adjoint
operators require domain care; the finite-dimensional
matrix rules are not automatically sufficient.}

\problem{3.5}{Adjoint rules and oscillator ladder operators}
Integration by parts shows, when the boundary
term vanishes, $(\partial_x)^\dagger=-\partial_x$.
Multiplication by $x$ satisfies $x^\dagger=x$,
whereas $(\ii I)^\dagger=-\ii I$. For operators
on a common compatible domain,
\begin{equation}
\boxed{(QR)^\dagger=R^\dagger Q^\dagger,\quad
(Q+R)^\dagger=Q^\dagger+R^\dagger,\quad
(cQ)^\dagger=c^*Q^\dagger}.
\end{equation}
These follow directly from applying the adjoint
definition twice: $\braket f{QRg}
=\braket{Q^\dagger f}{Rg}
=\braket{R^\dagger Q^\dagger f}g$.
The oscillator raising operator
$a_+=(m\omega x-\ii p)/\sqrt{2m\hbar\omega}$
therefore has adjoint
\begin{equation}
\boxed{a_+^\dagger=
\frac{m\omega x+\ii p}{\sqrt{2m\hbar\omega}}=a_-}.
\end{equation}

\problem{3.6}{Angular second derivative, eigenvalues and degeneracy}
On the periodic domain $f(\phi+2\pi)=f(\phi)$,
$f'(\phi+2\pi)=f'(\phi)$, integrate twice:
\begin{equation}
\int_0^{2\pi}f^*g''\dd\phi=
[f^*g'-f'^*g]_0^{2\pi}
+\int_0^{2\pi}f''{}^*g\dd\phi
=\int_0^{2\pi}f''{}^*g\dd\phi.
\end{equation}
Hence $Q=\partial_\phi^2$ is symmetric
(and self-adjoint on the usual periodic $H^2$ domain).
The normalized eigenfunctions are
\begin{equation}
\boxed{u_n(\phi)=\frac{\ee^{\ii n\phi}}{\sqrt{2\pi}}},
\quad\boxed{Q u_n=-n^2u_n},\qquad n\in\mathbb Z.
\end{equation}
The spectrum is $0,-1,-4,-9,\ldots$; $0$ is simple,
each $-n^2$ ($n\ge1$) is doubly degenerate,
with independent $\ee^{\pm\ii n\phi}$.

\problem{3.7}{Degenerate eigenspaces and orthogonal combinations}
If $Qf=qf$ and $Qg=qg$, linearity gives
$Q(\alpha f+\beta g)=q(\alpha f+\beta g)$.
For $Q=\partial_x^2$ on $(-1,1)$,
$f=\ee^x$ and $g=\ee^{-x}$ both have eigenvalue $1$.
Choose their even/odd combinations
\begin{equation}
u(x)=\cosh x=\frac{\ee^x+\ee^{-x}}2,\qquad
v(x)=\sinh x=\frac{\ee^x-\ee^{-x}}2.
\end{equation}
Both satisfy $u''=u$, $v''=v$, and
$\int_{-1}^1u(x)v(x)\dd x=0$
because the integrand is odd.
\qnote{These are orthogonal eigenfunctions of the
differential expression. A self-adjoint boundary
domain must be specified before interpreting the
degeneracy as a physical observable's spectrum.}

\problem{3.8}{Orthogonality of angular eigenfunctions}
For the angular momentum operator of Example 3.1,
$L_z=-\ii\hbar\partial_\phi$,
the periodic eigenfunctions
$u_n(\phi)=\ee^{\ii n\phi}/\sqrt{2\pi}$
obey $L_zu_n=n\hbar u_n$: the eigenvalues are real.
Their overlaps are
\begin{equation}
\braket{u_m}{u_n}=
\frac1{2\pi}\int_0^{2\pi}
\ee^{\ii(n-m)\phi}\dd\phi
=\delta_{mn}.
\end{equation}
For Problem 3.6, the same basis diagonalizes
$\partial_\phi^2$ with eigenvalues $-n^2$.
Distinct eigenvalues correspond to distinct
$|n|$, hence orthogonality follows from the
same integral. For $n$ and $-n$, orthogonality
still holds despite the degeneracy (for $n\ne0$).

\problem{3.9}{Discrete, continuous and mixed energy spectra}
A particle in a one-dimensional infinite well
has only the discrete energies
$E_n=n^2\pi^2\hbar^2/(2ma^2)$.
A potential step has only scattering states in
the customary no-bound-state configuration,
and hence a continuous energy spectrum.
An attractive delta well has one bound level
$E_b=-m\alpha^2/(2\hbar^2)$ and positive-energy
scattering states; its spectrum is mixed.
\qnote{Other examples may be used if the
potential and boundary conditions are stated:
the type of spectrum depends on the self-adjoint
Hamiltonian, not merely its differential formula.}

\problem{3.10}{Is the ground well state a momentum eigenstate?}
For $0<x<a$, the normalized ground state is
$\psi_1(x)=\sqrt{2/a}\sin(\pi x/a)$.
Its derivative gives
\begin{equation}
\hat p\psi_1=-\ii\hbar\frac{\pi}{a}
\sqrt{\frac2a}\cos(\pi x/a),
\end{equation}
which is not a constant multiple of $\psi_1$.
Therefore it is not a momentum eigenstate.
An energy eigenstate in a box is a standing wave,
not a single traveling momentum wave.
Although $\avg p=0$ and
$\avg{p^2}=\pi^2\hbar^2/a^2$, neither result
makes the state an eigenvector of $\hat p$.
\qnote{There is a subtle distinction between the
formal derivative and the domain of a self-adjoint
momentum operator on a finite interval,
explored explicitly in Problem 3.48.}

\section{Momentum representation and uncertainty (3.11–3.22)}
\problem{3.11}{Momentum distribution of the oscillator ground state}
Set $\ell=\sqrt{\hbar/(m\omega)}$. The normalized
position wavefunction is
$\psi_0(x)=(\pi\ell^2)^{-1/4}
\exp[-x^2/(2\ell^2)]$.
Fourier transformation with the $p$-normalization gives
\begin{align}
\Phi_0(p,t)&=\frac{\ee^{-\ii\omega t/2}}
{\sqrt{2\pi\hbar}}\int_{\mathbb R}
(\pi\ell^2)^{-1/4}
\ee^{-x^2/(2\ell^2)-\ii px/\hbar}\dd x\\
&=\boxed{(\pi m\hbar\omega)^{-1/4}
\ee^{-p^2/(2m\hbar\omega)}
\ee^{-\ii\omega t/2}}.
\end{align}
The classical momentum limit at $E_0=\hbar\omega/2$
is $p_c=\sqrt{2mE_0}=\sqrt{m\hbar\omega}$.
Therefore, with $u=p/\sqrt{m\hbar\omega}$,
\begin{equation}
P(|p|>p_c)=\frac2{\sqrt\pi}\int_1^\infty
\ee^{-u^2}\dd u=\operatorname{erfc}(1)
\simeq\boxed{0.16\quad\text{(two significant figures)}}.
\end{equation}
\qnote{The classically forbidden momentum tail
is the momentum-space counterpart of the
oscillator's position-space forbidden region.}

\problem{3.12}{Fourier coefficients of a free wave packet}
For the free-particle expansion
$\Psi(x,t)=
(2\pi)^{-1/2}\int\phi(k)
\ee^{\ii kx-\ii\hbar k^2t/(2m)}\dd k$,
the symmetric momentum Fourier transform yields
\begin{equation}
\Phi(p,t)=\frac1{\sqrt{2\pi\hbar}}
\int\ee^{-\ii px/\hbar}\Psi(x,t)\dd x.
\end{equation}
Insert the expansion and use
$\int_{\mathbb R}\ee^{\ii(k-p/\hbar)x}\dd x
=2\pi\delta(k-p/\hbar)$:
\begin{equation}
\boxed{\Phi(p,t)=\frac1{\sqrt\hbar}
\phi(p/\hbar)\ee^{-\ii p^2t/(2m\hbar)}}.
\end{equation}
Taking the squared modulus removes the phase:
\begin{equation}
\boxed{|\Phi(p,t)|^2=
\hbar^{-1}|\phi(p/\hbar)|^2},
\end{equation}
independent of time. This also follows
from the free Hamiltonian commuting with momentum.

\problem{3.13}{The position operator in momentum space}
Starting from the inverse transform,
$\Psi(x)=(2\pi\hbar)^{-1/2}\int
\ee^{\ii px/\hbar}\Phi(p)\dd p$,
differentiate the plane-wave factor:
$x\ee^{\ii px/\hbar}=
-\ii\hbar\,\partial_p\ee^{\ii px/\hbar}$.
Integration by parts (vanishing boundary term) then gives
\begin{equation}
x\Psi(x)=\frac1{\sqrt{2\pi\hbar}}
\int\ee^{\ii px/\hbar}
\left(\ii\hbar\,\partial_p\Phi(p)\right)\dd p.
\end{equation}
Since Fourier transformation is unitary,
\begin{equation}
\boxed{\avg x=\int_{\mathbb R}
\Phi^*(p)\,\ii\hbar\frac{\partial\Phi}{\partial p}\dd p}.
\end{equation}
Thus $\hat x=\ii\hbar\partial_p$
in the momentum representation, whereas
$\hat p$ acts by multiplication by $p$.

\problem{3.14}{Commutator algebra and functions of position}
Expand $\comm{A+B}C=(A+B)C-C(A+B)$ to obtain
$\boxed{\comm{A+B}C=\comm AC+\comm BC}$.
Similarly,
\begin{align}
\comm{AB}C
&=ABC-CAB\\
&=A(BC-CB)+(AC-CA)B\\
&=\boxed{A\comm BC+\comm AC\,B}.
\end{align}
The canonical commutator gives
$\comm{x^n}p=
x^{n-1}\comm xp+
\comm{x^{n-1}}p\,x$.
By induction,
\begin{equation}
\boxed{\comm{x^n}p=\ii n\hbar\,x^{n-1}}.
\end{equation}
If $f(x)=\sum_{n=0}^\infty a_nx^n$
and formal termwise differentiation is valid
on a common dense domain,
\begin{equation}
\boxed{\comm{f(x)}p=
\sum_{n=1}^\infty a_n\ii n\hbar x^{n-1}
=\ii\hbar f'(x)}.
\end{equation}
Direct verification:
$f(-\ii\hbar\psi')+
\ii\hbar(f\psi)'=\ii\hbar f'\psi$.

\problem{3.15}{Position–energy uncertainty relation}
Apply Robertson's inequality to $A=x$ and $B=H$:
$\sigma_x\sigma_H\ge
\frac12|\avg{\comm{x}{H}}|$.
For $H=p^2/(2m)+V(x)$,
$\comm{x}{V(x)}=0$ and
\begin{equation}
\comm{x}{p^2}=\comm xp\,p+
p\comm xp=2\ii\hbar p.
\end{equation}
Consequently
\begin{equation}
\boxed{\sigma_x\sigma_H\ge
\frac{\hbar}{2m}|\avg p|}.
\end{equation}
In a stationary energy eigenstate,
$\sigma_H=0$, but also
$\dd\avg x/\dd t=\avg p/m=0$
(for a normalizable state with no boundary
current). Therefore $\avg p=0$
and the inequality becomes $0\ge0$.

\problem{3.16}{Noncommutativity forbids a complete common eigenbasis}
Assume two linear operators $A,B$ have
a complete basis of simultaneous eigenvectors
$u_n$, with $Au_n=a_nu_n$ and $Bu_n=b_nu_n$.
For every basis vector,
\begin{equation}
[A,B]u_n=(a_nb_n-b_na_n)u_n=0.
\end{equation}
By linearity, $[A,B]$ annihilates every
finite linear combination of basis vectors.
On a finite-dimensional space this gives
$[A,B]=0$ on the whole space. For
unbounded operators in an infinite-dimensional
Hilbert space, the same conclusion holds
on a common invariant dense span, subject
to the required domains. The contrapositive is
that a genuinely nonzero commutator excludes
a complete common eigenbasis.

\problem{3.17}{All normalized minimum-uncertainty position states}
Equality in Cauchy--Schwarz requires
$(p-p_0)\Psi=\ii c(x-x_0)\Psi$
for an appropriate real nonzero constant $c$.
In the position representation,
\begin{equation}
-\ii\hbar\Psi'(x)-p_0\Psi(x)
=\ii c(x-x_0)\Psi(x).
\end{equation}
Divide by $\Psi$ and integrate:
\begin{equation}
\frac{\Psi'}{\Psi}=\frac{\ii p_0}{\hbar}
-\frac{c}{\hbar}(x-x_0),
\quad
\Psi(x)=A\exp\left[
-\frac{c(x-x_0)^2}{2\hbar}+
\frac{\ii p_0x}{\hbar}\right].
\end{equation}
Normalizability requires $c>0$. Writing
$c=\hbar/(2\sigma_x^2)$ and normalizing yields
\begin{equation}
\boxed{\Psi(x)=
(2\pi\sigma_x^2)^{-1/4}
\exp\left[-\frac{(x-x_0)^2}{4\sigma_x^2}
+\frac{\ii p_0x}{\hbar}\right]}.
\end{equation}
Its distribution has
$\avg x=x_0$, $\sigma_x$ as specified,
$\avg p=p_0$, $\sigma_p=\hbar/(2\sigma_x)$,
hence $\sigma_x\sigma_p=\hbar/2$.
\qnote{A physically irrelevant global phase
may multiply the solution.}

\problem{3.18}{Time derivatives of four expectation values}
For time-independent $H$ and any operator $Q$,
\begin{equation}
\frac{\dd}{\dd t}\avg Q=
\frac{\ii}{\hbar}\avg{\comm H Q}
+\left\langle\frac{\partial Q}{\partial t}\right\rangle.
\end{equation}
(a) $Q=I$: $\comm HI=0$ gives
$\boxed{\dd\avg I/\dd t=0}$, norm conservation.
(b) $Q=H$: $\comm HH=0$ gives
$\boxed{\dd\avg H/\dd t=0}$, energy conservation
when $\partial H/\partial t=0$.
(c) $Q=x$: $\comm Hx=-\ii\hbar p/m$ gives
\begin{equation}
\boxed{\dd\avg x/\dd t=\avg p/m}.
\end{equation}
(d) $Q=p$: $\comm Hp=\ii\hbar V'(x)$ gives
\begin{equation}
\boxed{\dd\avg p/\dd t=-\avg{V'(x)}}.
\end{equation}
These are the Ehrenfest equations.

\problem{3.19}{Free and harmonic-oscillator motion of the mean}
For a free particle, $V=0$, hence
$\dd\avg p/\dd t=0$. Integrating twice gives
\begin{equation}
\boxed{\avg x(t)=\avg x(0)+
\frac{\avg p(0)}m t}.
\end{equation}
For $V(x)=m\omega^2x^2/2$, Ehrenfest's
equations imply
$m\dd^2\avg x/\dd t^2=-m\omega^2\avg x$.
With $x_0=\avg x(0)$ and $p_0=\avg p(0)$,
\begin{equation}
\boxed{\avg x(t)=x_0\cos\omega t+
\frac{p_0}{m\omega}\sin\omega t}.
\end{equation}
Both results hold for any normalized state
with the appropriate finite expectations.

\problem{3.20}{Energy–time inequality for a two-level box superposition}
Use the normalized state from Problem 2.5,
$(\phi_1\ee^{-\ii E_1t/\hbar}
+\phi_2\ee^{-\ii E_2t/\hbar})/\sqrt2$,
$E_n=n^2\hbar\Omega$,
$\Omega=\pi^2\hbar/(2ma^2)$.
Its energy uncertainty is
\begin{equation}
\sigma_H=\frac{|E_2-E_1|}{2}
=\boxed{\frac{3\hbar\Omega}{2}}.
\end{equation}
By integrating sine products,
$\avg x=a/2-\frac{16a}{9\pi^2}\cos(3\Omega t)$,
$\avg{x^2}=a^2(\frac13-\frac{5}{16\pi^2})
-\frac{16a^2}{9\pi^2}\cos(3\Omega t)$, so
\begin{equation}
\boxed{\sigma_x^2(t)=a^2
\left[\frac1{12}-\frac{5}{16\pi^2}
-\frac{256}{81\pi^4}\cos^2(3\Omega t)\right]}.
\end{equation}
Furthermore
\begin{equation}
\left|\frac{\dd\avg x}{\dd t}\right|=
\frac{16a\Omega}{3\pi^2}
|\sin(3\Omega t)|.
\end{equation}
With $\Delta t=\sigma_x/|\dd\avg x/\dd t|$
(where the denominator is nonzero),
\begin{equation}
\sigma_H\Delta t=
\frac{9\pi^2\hbar}{32a}
\frac{\sigma_x(t)}{|\sin(3\Omega t)|}
\ge\frac{\hbar}{2}.
\end{equation}
To check the last inequality algebraically,
square it and use
$\sigma_x^2/a^2\ge
(256/(81\pi^4))\sin^2(3\Omega t)$;
the difference is
$1/12-5/(16\pi^2)-256/(81\pi^4)>0$.
At instants when the mean position is momentarily
stationary, $\Delta t$ is infinite or undefined.

\problem{3.21}{Energy–time inequality for a traveling Gaussian}
For the freely moving Gaussian of Problem 2.42,
$\Psi(x,0)\propto\ee^{-ax^2+\ii \ell x}$,
where $a>0$, the momentum distribution is
Gaussian with mean $p_0=\hbar\ell$ and
variance $s_p^2=a\hbar^2$.
For a normal random variable,
$\mathrm{Var}(p^2)=4p_0^2s_p^2+2s_p^4$;
hence, with $H=p^2/(2m)$,
\begin{equation}
\boxed{\sigma_H=
\frac{\hbar^2}{m}\sqrt{a\ell^2+\frac{a^2}{2}}}.
\end{equation}
Writing $s=2\hbar at/m$, the wavepacket width
and mean velocity are
\begin{equation}
\boxed{\sigma_x(t)=\frac{\sqrt{1+s^2}}{2\sqrt a}},
\quad
\boxed{\frac{\dd\avg x}{\dd t}=\frac{\hbar\ell}{m}}.
\end{equation}
For $\ell\ne0$,
\begin{equation}
\sigma_H\Delta t=
\frac{\hbar}{2}\sqrt{1+\frac{a}{2\ell^2}}
\sqrt{1+s^2}\ge\frac{\hbar}{2}.
\end{equation}
If $\ell=0$, the mean does not move and this
particular characteristic-time definition
is not informative.

\problem{3.22}{Energy–time and position–energy uncertainty coincide}
For $Q=x$ the Mandelstam–Tamm timescale is
$\Delta t=\sigma_x/|\dd\avg x/\dd t|$.
Ehrenfest gives $\dd\avg x/\dd t=\avg p/m$,
so the energy–time inequality
$\sigma_H\Delta t\ge\hbar/2$ becomes
\begin{equation}
\frac{\sigma_H\sigma_x m}{|\avg p|}
\ge\frac\hbar2
\quad\Longleftrightarrow\quad
\boxed{\sigma_x\sigma_H\ge
\frac{\hbar|\avg p|}{2m}},
\end{equation}
which is exactly Problem 3.15.
When $\avg p=0$ the intermediate division
is not permitted; the direct commutator form
remains valid.

\section{Dirac notation, matrices and transformations (3.23–3.30)}
\problem{3.23}{Projection operators and their spectrum}
For a normalized ket $\ket\alpha$ define
$P=\ket\alpha\bra\alpha$. Then
\begin{equation}
P^2=\ket\alpha\braket\alpha\alpha\bra\alpha
=\ket\alpha\bra\alpha=\boxed P.
\end{equation}
Any eigenvalue $p$ obeys $p^2=p$, hence
$p=0$ or $1$. The vector $\ket\alpha$
is a $p=1$ eigenvector, and every vector
orthogonal to it is in the $p=0$
eigenspace. A projection onto a
higher-dimensional subspace has the same
eigenvalues, with the range and kernel
as the corresponding eigenspaces.

\problem{3.24}{Hermiticity of an operator's matrix}
Let $\{\ket n\}$ be an orthonormal basis and
$Q_{mn}=\bra m Q\ket n$.
If $Q=Q^\dagger$,
\begin{align}
Q_{mn}&=\braket m{Qn}
=\braket{Qm}n\\
&=\braket n{Qm}^*
=\boxed{Q_{nm}^*}.
\end{align}
Thus the matrix $Q$ satisfies
$\boxed{Q=Q^\dagger}$, i.e.\ its
conjugate transpose equals itself.
In particular all diagonal entries are real.

\problem{3.25}{Diagonalizing a two-level Hamiltonian}
In the orthonormal basis $\ket1,\ket2$,
the stated operator has the matrix
\begin{equation}
H=\varepsilon
\begin{pmatrix}1&1\\1&-1\end{pmatrix}.
\end{equation}
Its characteristic equation is
\begin{equation}
\det(H-EI)=E^2-2\varepsilon^2=0,
\qquad\boxed{E_\pm=\pm\sqrt2\,\varepsilon}.
\end{equation}
For $E_+$ the first row gives
$v_2=(\sqrt2-1)v_1$; for $E_-$,
$v_2=-(\sqrt2+1)v_1$.
A normalized eigenbasis is therefore
\begin{align}
\boxed{\ket{E_+}}&=
\frac{\ket1+(\sqrt2-1)\ket2}
{\sqrt{1+(\sqrt2-1)^2}},\\
\boxed{\ket{E_-}}&=
\frac{\ket1-(\sqrt2+1)\ket2}
{\sqrt{1+(\sqrt2+1)^2}}.
\end{align}
The vectors are orthogonal because
$1+(\sqrt2-1)(-\sqrt2-1)=0$.

\problem{3.26}{Complex kets, dual vectors and outer products}
With basis coordinates
$\ket\alpha=(\ii,-2,-\ii)^{\mathsf T}$,
$\ket\beta=(\ii,0,2)^{\mathsf T}$,
taking conjugate transposes gives
\begin{equation}
\boxed{\bra\alpha=(-\ii,-2,\ii)},\qquad
\boxed{\bra\beta=(-\ii,0,2)}.
\end{equation}
The two inner products are
\begin{equation}
\boxed{\braket\alpha\beta=
(-\ii)\ii+\ii(2)=1+2\ii},\qquad
\boxed{\braket\beta\alpha=1-2\ii}.
\end{equation}
The outer product $A=\ket\alpha\bra\beta$ is
\begin{equation}
\boxed{A=
\begin{pmatrix}
1&0&2\ii\\
2\ii&0&-4\\
-1&0&-2\ii
\end{pmatrix}}.
\end{equation}
For example $A_{13}=\ii\cdot2=2\ii$ and
$A_{31}=(-\ii)(-\ii)=-1$.
Since $A_{13}\ne A_{31}^*$,
$\boxed{A\ne A^\dagger}$.

\problem{3.27}{Spectral decomposition and functions of an operator}
Let $Q\ket{e_n}=q_n\ket{e_n}$, where
$\{\ket{e_n}\}$ is a complete orthonormal
basis. Completeness gives
$\sum_n\ket{e_n}\bra{e_n}=I$. On an arbitrary
vector $\ket\alpha=\sum_n\alpha_n\ket{e_n}$,
\begin{equation}
Q\ket\alpha=
\sum_nq_n\alpha_n\ket{e_n}=
\left(\sum_nq_n\ket{e_n}\bra{e_n}\right)\ket\alpha.
\end{equation}
Therefore
\begin{equation}
\boxed{Q=\sum_nq_n\ket{e_n}\bra{e_n}}.
\end{equation}
For a function $f$ defined on the spectrum,
$f(Q)=\sum_nf(q_n)\ket{e_n}\bra{e_n}$.
For the exponential, its power series gives
\begin{align}
\ee^Q&=\sum_{k=0}^\infty Q^k/k!\\
&=\sum_n\sum_{k=0}^\infty
q_n^k/k!\ket{e_n}\bra{e_n}\\
&=\boxed{\sum_n\ee^{q_n}\ket{e_n}\bra{e_n}}.
\end{align}
\qnote{For unbounded operators, convergence and
the domain of each function of $Q$ are
specified by the spectral theorem.}

\problem{3.28}{Functions of the differentiation operator}
For $D=\dd/\dd x$, apply
$\sin D=\sum_{j\ge0}(-1)^jD^{2j+1}/(2j+1)!$
to $x^5$. Higher than fifth derivatives vanish:
\begin{equation}
\boxed{(\sin D)x^5=
5x^4-\frac{60x^2}{6}+\frac{120}{120}
=5x^4-10x^2+1}.
\end{equation}
To evaluate $(1-D/2)^{-1}\cos x$, solve
$(1-D/2)f=\cos x$. Set
$f=a\cos x+b\sin x$.
Since $D\cos x=-\sin x$ and
$D\sin x=\cos x$,
\begin{equation}
(a-b/2)\cos x+(b+a/2)\sin x=\cos x.
\end{equation}
The equations $a-b/2=1$, $b+a/2=0$
give $a=4/5$ and $b=-2/5$, hence
\begin{equation}
\boxed{(1-D/2)^{-1}\cos x=
\frac45\cos x-\frac25\sin x}.
\end{equation}
\qnote{The inverse is interpreted on the
two-dimensional invariant span of
$\cos x$ and $\sin x$, where it is
unambiguously invertible.}

\problem{3.29}{A central-commutator Baker–Campbell–Hausdorff formula}
Assume $C=[A,B]$ and $[A,C]=[B,C]=0$.
Induction from
$[A^{n+1},B]=A[A^n,B]+[A,B]A^n$
gives
\begin{equation}
\boxed{[A^n,B]=nA^{n-1}C}.
\end{equation}
Expanding $\ee^{\lambda A}$ into a series,
\begin{align}
[\ee^{\lambda A},B]
&=\sum_{n=1}^\infty\frac{\lambda^n}{n!}
[A^n,B]\\
&=\lambda\ee^{\lambda A}C.
\end{align}
To establish the BCH identity, define
$F(t)=\ee^{t(A+B)}$ and
$G(t)=\ee^{tA}\ee^{tB}\ee^{-t^2C/2}$.
Using $[A,\ee^{tB}]=tC\ee^{tB}$,
or equivalently
$\ee^{tA}B=(B+tC)\ee^{tA}$,
differentiate $G$:
\begin{equation}
G'=(A+B+tC-tC)G=(A+B)G.
\end{equation}
Since $G(0)=I=F(0)$ and
$F'=(A+B)F$, uniqueness gives
$F(t)=G(t)$ and at $t=1$,
\begin{equation}
\boxed{\ee^{A+B}=\ee^A\ee^B\ee^{-C/2}}.
\end{equation}
\qnote{For unbounded operators, this algebra
is formal unless the exponential operators
and their domains satisfy the necessary
analytic conditions. The Weyl relations
provide a rigorous form for $x,p$.}

\problem{3.30}{From position-space to energy-space amplitudes}
Let $\ket n$ be a discrete orthonormal
energy eigenbasis, with position representation
$\psi_n(x)=\braket x n$. For a state $\ket{S(t)}$,
insert the position resolution of the identity:
\begin{align}
c_n(t)&=\braket n{S(t)}
=\bra n\left(\int_{\mathbb R}
\ket x\bra x\dd x\right)\ket{S(t)}\\
&=\int_{\mathbb R}
\braket n x\,\braket x{S(t)}\,\dd x\\
&=\boxed{\int_{\mathbb R}
\psi_n^*(x)\Psi(x,t)\dd x}.
\end{align}
For a time-independent Hamiltonian,
\begin{equation}
\boxed{c_n(t)=c_n(0)\ee^{-\ii E_nt/\hbar}},
\quad\boxed{\Psi(x,t)=\sum_nc_n(t)\psi_n(x)}.
\end{equation}
This is the energy-basis transformation
of the same Hilbert-space vector.

\section{Further formalism and quantum statistics (3.31–3.40)}
\problem{3.31}{Orthonormal Legendre polynomials from Gram–Schmidt}
On $[-1,1]$ take the inner product
$\braket fg=\int_{-1}^1 f(x)g(x)\dd x$.
Start with $u_0=1$, $\|u_0\|=\sqrt2$.
Since $\int x\dd x=0$, $u_1=x$,
$\|u_1\|=\sqrt{2/3}$.
For $x^2$, subtract the projection onto $1$:
\begin{equation}
u_2=x^2-\frac{\int_{-1}^1x^2\dd x}
{\int_{-1}^1 1\dd x}=x^2-\frac13.
\end{equation}
It is automatically orthogonal to $x$ by parity,
and $\|u_2\|^2=8/45$.
For $x^3$, subtract the projection onto $x$:
\begin{equation}
u_3=x^3-\frac{\int x^4\dd x}
{\int x^2\dd x}\,x=x^3-\frac35x.
\end{equation}
The other projections vanish by parity or
orthogonality, and $\|u_3\|^2=8/175$.
The orthonormal basis is therefore
\begin{align}
\boxed{e_0(x)}&=\frac1{\sqrt2},&
\boxed{e_1(x)}&=\sqrt{\frac32}\,x,\\
\boxed{e_2(x)}&=\sqrt{\frac52}\,
\frac{3x^2-1}{2},&
\boxed{e_3(x)}&=\sqrt{\frac72}\,
\frac{5x^3-3x}{2}.
\end{align}
These are $\sqrt{(2n+1)/2}\,P_n(x)$
for Legendre polynomials normalized by $P_n(1)=1$.

\problem{3.32}{Anti-Hermitian operators and their commutators}
Let $Q^\dagger=-Q$.
Then
$\avg Q^*=\bra\psi Q^\dagger\ket\psi=-\avg Q$,
so $\avg Q$ is purely imaginary (or zero).
For $Q\ket u=\lambda\ket u$ with $\|u\|=1$,
$\lambda=\avg Q$ is likewise imaginary.
For eigenvectors $u,v$ with eigenvalues
$\lambda,\mu$, respectively,
\begin{equation}
\braket{Qu}v=\lambda^*\braket uv
=-\braket u{Qv}=-\mu\braket uv.
\end{equation}
Hence $(\lambda^*+\mu)\braket uv=0$.
Since $\lambda^*=-\lambda$,
distinct eigenvalues imply
$\boxed{\braket uv=0}$.
For Hermitian $A,B$,
\begin{equation}
[A,B]^\dagger=(AB-BA)^\dagger=BA-AB=-[A,B].
\end{equation}
The same conclusion holds if both operators
are anti-Hermitian: their commutator is
again anti-Hermitian.
Finally every operator $Q$ decomposes as
\begin{equation}
\boxed{Q=A+B,\quad A=\frac{Q+Q^\dagger}{2},
\quad B=\frac{Q-Q^\dagger}{2}},
\end{equation}
with $A^\dagger=A$ and $B^\dagger=-B$.

\problem{3.33}{Three sequential projective measurements}
After measuring $A$ and obtaining $a_1$,
the state is $\psi_1=(3\phi_1+4\phi_2)/5$.
A subsequent measurement of $B$ yields
$b_1$ with probability $9/25$ or $b_2$
with probability $16/25$.
Conditional on $b_1$, the state becomes
$\phi_1$ and the next probability of
$a_1$ is $|\braket{\psi_1}{\phi_1}|^2=9/25$;
conditional on $b_2$, that probability is
$16/25$.
If the $B$ outcome is not disclosed,
the law of total probability gives
\begin{equation}
\boxed{P(a_1\text{ at the third measurement})=
\frac9{25}\frac9{25}+
\frac{16}{25}\frac{16}{25}
=\frac{337}{625}}.
\end{equation}
The result differs from the probability one
without an intervening nonselective $B$ measurement,
because that measurement destroys coherence
in the $B$ basis.

\problem{3.34}{Infinite-well energy eigenstates in momentum space}
For $0<x<a$ write
$\psi_n(x)=\sqrt{2/a}\sin(k_nx)$,
$k_n=n\pi/a$, and zero otherwise.
The Fourier transform integral is
\begin{align}
\Phi_n(p,t)&=\frac{\ee^{-\ii E_nt/\hbar}}
{\sqrt{2\pi\hbar}}\sqrt{\frac2a}
\int_0^a\sin(k_nx)\ee^{-\ii px/\hbar}\dd x\\
&=\boxed{
\frac{\ee^{-\ii E_nt/\hbar}}{\sqrt{\pi a\hbar}}
\frac{k_n[1-(-1)^n\ee^{-\ii pa/\hbar}]}
{k_n^2-(p/\hbar)^2}}.
\end{align}
Apparent poles at $p=\pm\hbar k_n$ are
removable: the numerator vanishes there too.
Writing $q=p/\hbar$, the probability density is
\begin{equation}
\boxed{|\Phi_n(p,t)|^2=
\frac{2k_n^2[1-(-1)^n\cos(qa)]}
{\pi a\hbar(k_n^2-q^2)^2}}.
\end{equation}
It is symmetric in $p$, hence $\avg p=0$.
For large $n$, the dominant momentum peaks
are near $p=\pm n\pi\hbar/a$, reflecting the
two traveling-wave constituents of the standing
energy eigenstate. Parseval's identity for
Fourier transforms gives
\begin{equation}
\avg{p^2}=
\hbar^2\int_{\mathbb R}|\psi_n'(x)|^2\dd x
=\boxed{\left(\frac{n\pi\hbar}{a}\right)^2},
\end{equation}
in agreement with Problem 2.4.
\qnote{The PDF's numerical plots use a
dimensionless momentum $pa/\hbar$ and
$n=1,2,5,10$.}

\problem{3.35}{Finite wave train and the cost of sharp edges}
Set $L=n\lambda$ with positive integer $n$,
$k_0=2\pi/\lambda$, and
$\Psi(x,0)=\ee^{\ii k_0x}/\sqrt{2L}$
on $-L<x<L$, zero elsewhere.
Its Fourier transform is
\begin{align}
\Phi(p,0)
&=\frac1{\sqrt{4\pi L\hbar}}
\int_{-L}^L\ee^{-\ii(p/\hbar-k_0)x}\dd x\\
&=\boxed{\frac1{\sqrt{\pi L\hbar}}
\frac{\sin[(p/\hbar-k_0)L]}{p/\hbar-k_0}}.
\end{align}
The position density is a rectangle of height
$1/(2L)$ and width $w_x=2L$.
The momentum density is a squared-sinc distribution
centered at $p_0=\hbar k_0$.
Its first zeros are at $p=p_0\pm\pi\hbar/L$,
so the main-lobe width is
\begin{equation}
\boxed{w_p=\frac{2\pi\hbar}{L}},\qquad
\boxed{w_xw_p=4\pi\hbar}.
\end{equation}
As $n\to\infty$, $w_x\to\infty$ and
$w_p\to0$. These widths are only
rough spread estimates, not standard deviations.
Indeed $|\Phi(p)|^2$ decays as $p^{-2}$
with nonvanishing oscillations. Therefore
$\int p^2|\Phi(p)|^2\dd p$ diverges and
$\sigma_p=\infty$. The sharp cutoff
creates jumps in $\Psi$ and invalidates a
finite kinetic-energy variance.
\qnote{One must not claim a finite Heisenberg
standard-deviation product based on the
sinc peak widths.}

\problem{3.36}{Lorentzian-amplitude wavefunction and its Fourier dual}
Take $\Psi(x,0)=A/(x^2+a^2)$ with $a>0$.
Use
\begin{equation}
\int_{-\infty}^\infty\frac{\dd x}{(x^2+a^2)^2}
=\frac{\pi}{2a^3},
\quad
\int_{-\infty}^\infty\frac{x^2\dd x}{(x^2+a^2)^2}
=\frac{\pi}{2a}.
\end{equation}
Normalization gives
$\boxed{A=\sqrt{2a^3/\pi}}$.
Parity yields $\avg x=0$, and the second
integral gives $\avg{x^2}=a^2$,
so $\boxed{\sigma_x=a}$.
The standard Fourier integral
$\int\ee^{-\ii kx}/(x^2+a^2)\dd x
=(\pi/a)\ee^{-a|k|}$ gives
\begin{equation}
\boxed{\Phi(p,0)=
\sqrt{\frac a\hbar}\ee^{-a|p|/\hbar}}.
\end{equation}
Check normalization:
$\int|\Phi|^2\dd p=
(2a/\hbar)\int_0^\infty\ee^{-2ap/\hbar}\dd p=1$.
Again $\avg p=0$, and
\begin{equation}
\avg{p^2}=\frac{2a}{\hbar}
\int_0^\infty p^2\ee^{-2ap/\hbar}\dd p
=\frac{\hbar^2}{2a^2}.
\end{equation}
Thus
\begin{equation}
\boxed{\sigma_p=\frac{\hbar}{\sqrt2a}},
\qquad
\boxed{\sigma_x\sigma_p=\frac{\hbar}{\sqrt2}
>\frac{\hbar}{2}}.
\end{equation}

\problem{3.37}{General virial theorem and the harmonic oscillator}
For $H=p^2/(2m)+V(x)$, apply
Ehrenfest's identity to the operator $xp$
(on an appropriate domain):
\begin{equation}
\frac{\dd}{\dd t}\avg{xp}=
\frac{\ii}{\hbar}\avg{[H,xp]}.
\end{equation}
Use $[H,xp]=[H,x]p+x[H,p]$
and the commutators
$[H,x]=-\ii\hbar p/m$,
$[H,p]=\ii\hbar V'(x)$:
\begin{equation}
\boxed{\frac{\dd}{\dd t}\avg{xp}
=2\avg T-\avg{xV'(x)}}.
\end{equation}
For a stationary energy eigenstate, the
expectation of a time-independent observable
is constant, hence the left side vanishes:
\begin{equation}
\boxed{2\avg T=\avg{xV'(x)}}.
\end{equation}
For $V=m\omega^2x^2/2$,
$xV'=m\omega^2x^2=2V$,
so $\boxed{\avg T=\avg V=E_n/2}$,
consistent with oscillator energy moments.

\problem{3.38}{Time to become orthogonal for two energy levels}
At $t=0$ let
$\ket{\Psi(0)}=(\ket1+\ket2)/\sqrt2$
with orthonormal energy states and
$E_1\ne E_2$.
Time evolution gives
\begin{equation}
\ket{\Psi(t)}=\frac1{\sqrt2}
(\ee^{-\ii E_1t/\hbar}\ket1+
\ee^{-\ii E_2t/\hbar}\ket2).
\end{equation}
Its overlap with the initial state is
\begin{align}
\braket{\Psi(0)}{\Psi(t)}
&=\frac12(\ee^{-\ii E_1t/\hbar}
+\ee^{-\ii E_2t/\hbar})\\
&=\ee^{-\ii(E_1+E_2)t/(2\hbar)}
\cos\left(\frac{(E_2-E_1)t}{2\hbar}\right).
\end{align}
The first orthogonality time is
$\boxed{\tau=\pi\hbar/|E_2-E_1|}$.
The energy standard deviation is
$\boxed{\sigma_H=|E_2-E_1|/2}$.
With the specified $\Delta t=\tau/\pi$,
\begin{equation}
\boxed{\sigma_H\Delta t=
\frac{|E_2-E_1|}{2}\,
\frac{\hbar}{|E_2-E_1|}
=\frac{\hbar}{2}}.
\end{equation}
The case of equal energy does not
orthogonalize and has $\sigma_H=0$.

\problem{3.39}{Oscillator matrices for position, momentum and Hamiltonian}
Let $a\ket n=\sqrt n\ket{n-1}$,
$a^\dagger\ket n=\sqrt{n+1}\ket{n+1}$.
With $\ell_x=\sqrt{\hbar/(2m\omega)}$,
$\ell_p=\sqrt{m\hbar\omega/2}$,
\begin{align}
x&=\ell_x(a+a^\dagger),&
p&=\ii\ell_p(a^\dagger-a).
\end{align}
Therefore, for nonnegative indices $j,k$,
\begin{align}
\boxed{X_{jk}}&=
\ell_x(\sqrt{k}\,\delta_{j,k-1}
+\sqrt{k+1}\,\delta_{j,k+1}),\\
\boxed{P_{jk}}&=
\ii\ell_p(\sqrt{k+1}\,\delta_{j,k+1}
-\sqrt{k}\,\delta_{j,k-1}).
\end{align}
Thus $X$ is real symmetric and tridiagonal,
and $P$ is imaginary Hermitian and tridiagonal.
Squaring in the \emph{complete infinite basis},
\begin{equation}
\frac{P^2}{2m}+\frac12m\omega^2X^2
=\frac{\hbar\omega}{4}
[(a^\dagger-a)^2(-1)+(a+a^\dagger)^2]
=\frac{\hbar\omega}{2}(aa^\dagger+a^\dagger a).
\end{equation}
With $aa^\dagger=N+1$ and $a^\dagger a=N$,
\begin{equation}
\boxed{H_{jk}=\hbar\omega(k+\tfrac12)\delta_{jk}}.
\end{equation}
\qnote{Truncating $X$ and $P$ before squaring
can corrupt their topmost diagonal level.}

\problem{3.40}{Classical-frequency motion for any oscillator state}
Let
$\ket{\Psi(t)}=\sum_{n=0}^\infty
c_n\ee^{-\ii\omega(n+1/2)t}\ket n$
with $\sum|c_n|^2=1$.
The oscillator ladder matrix elements give
\begin{equation}
\avg{a}(t)=\ee^{-\ii\omega t}
\sum_{n=0}^\infty\sqrt{n+1}\,c_n^*c_{n+1}.
\end{equation}
Set $S=\sum\sqrt{n+1}\,c_n^*c_{n+1}$
and write
$S=|S|\ee^{\ii\phi}$.
Since
$x=\sqrt{\hbar/(2m\omega)}(a+a^\dagger)$,
\begin{equation}
\boxed{\avg x(t)=C\cos(\omega t-\phi)},
\quad
\boxed{C=\sqrt{\frac{2\hbar}{m\omega}}|S|}.
\end{equation}
Equivalently
$C\ee^{-\ii\phi}=
\sqrt{2\hbar/(m\omega)}
\sum\sqrt{n+1}\,c_{n+1}^*c_n$.
For Problem 2.40, the real coefficients are
$c_0=3/5$, $c_1=-2\sqrt2/5$,
$c_2=2\sqrt2/5$, with all others zero.
Then
\begin{equation}
S=c_0c_1+\sqrt2c_1c_2
=-\frac{14\sqrt2}{25},
\end{equation}
and
\begin{equation}
\boxed{C=\frac{28}{25}
\sqrt{\frac{\hbar}{m\omega}},\quad \phi=\pi}.
\end{equation}
Thus every oscillator state with finite
$\avg x$ follows the \emph{classical
frequency} $\omega$, even if its
probability density is highly nonclassical.

\section{Coherent states, supersymmetry and operator domains (3.41–3.49)}
\problem{3.41}{Largest momentum expectation for two oscillator levels}
Equal probabilities for the first two energies imply
the general normalized state
\begin{equation}
\ket{\Psi(0)}=\frac1{\sqrt2}
(\ket0+\ee^{\ii\theta}\ket1),
\end{equation}
where a global phase has been removed.
Because $p=\ii\sqrt{m\hbar\omega/2}(a^\dagger-a)$
and $\avg a=\ee^{\ii\theta}/2$,
\begin{equation}
\boxed{\avg p(0)=
\sqrt{\frac{m\hbar\omega}{2}}\sin\theta}.
\end{equation}
The largest \emph{positive} value is
$\boxed{\sqrt{m\hbar\omega/2}}$,
achieved at $\theta=\pi/2$.
Including time dependence,
\begin{equation}
\boxed{\Psi(x,t)=\frac1{\sqrt2}
\left[\psi_0(x)\ee^{-\ii\omega t/2}+
\ii\psi_1(x)\ee^{-3\ii\omega t/2}\right]}.
\end{equation}
The momentum expectation later evolves
as $\avg p(t)=\sqrt{m\hbar\omega/2}\cos\omega t$.

\problem{3.42}{Coherent states, their normalization and time evolution}
Let $a\ket\alpha=\alpha\ket\alpha$,
where $\alpha$ is arbitrary complex.
Taking the adjoint gives
$\bra\alpha a^\dagger=\alpha^*\bra\alpha$.
For $x=\ell_x(a+a^\dagger)$,
$p=\ii\ell_p(a^\dagger-a)$,
with $\ell_x=\sqrt{\hbar/(2m\omega)}$ and
$\ell_p=\sqrt{m\hbar\omega/2}$,
\begin{align}
\boxed{\avg x}&=\boxed{\ell_x(\alpha+\alpha^*)},&
\boxed{\avg p}&=\boxed{\ii\ell_p(\alpha^*-\alpha)},\\
\boxed{\avg{x^2}}&=\boxed{\avg x^2+\ell_x^2},&
\boxed{\avg{p^2}}&=\boxed{\avg p^2+\ell_p^2}.
\end{align}
Indeed $\avg{a^\dagger a}=|\alpha|^2$,
$\avg{aa^\dagger}=|\alpha|^2+1$,
$\avg{a^2}=\alpha^2$ and
$\avg{(a^\dagger)^2}=(\alpha^*)^2$.
It follows that
\begin{equation}
\boxed{\sigma_x=\sqrt{\frac{\hbar}{2m\omega}}},
\quad
\boxed{\sigma_p=\sqrt{\frac{m\hbar\omega}{2}}},
\quad
\boxed{\sigma_x\sigma_p=\frac\hbar2}.
\end{equation}
Expand $\ket\alpha=\sum_{n=0}^\infty c_n\ket n$.
Comparing coefficients in
$a\ket\alpha=\alpha\ket\alpha$ yields
$c_{n+1}\sqrt{n+1}=\alpha c_n$, hence
\begin{equation}
c_n=\frac{\alpha^n}{\sqrt{n!}}c_0.
\end{equation}
Normalization gives
\begin{equation}
1=|c_0|^2\sum_{n=0}^\infty
\frac{|\alpha|^{2n}}{n!}
=|c_0|^2\ee^{|\alpha|^2}
\quad\Rightarrow\quad
\boxed{c_0=\ee^{-|\alpha|^2/2}}
\end{equation}
up to an overall phase.
Under oscillator evolution,
\begin{equation}
\ket{\Psi(t)}=
\ee^{-\ii\omega t/2}
\ket{\alpha\ee^{-\ii\omega t}},
\quad
\boxed{\alpha(t)=\alpha\ee^{-\ii\omega t}}.
\end{equation}
Thus the state remains coherent, with fixed
uncertainties. If
$\alpha=C\sqrt{m\omega/(2\hbar)}
\ee^{\ii\phi}$, then
\begin{equation}
\boxed{\avg{x}(t)=C\cos(\omega t-\phi)},
\qquad
\boxed{\sigma_x(t)=\sqrt{\frac{\hbar}{2m\omega}}}.
\end{equation}
The ground state $\ket0$ is coherent with
$\boxed{\alpha=0}$.

\problem{3.43}{Robertson–Schrödinger uncertainty bound}
For Hermitian $A,B$, set
$\delta A=A-\avg A\,I$,
$\delta B=B-\avg B\,I$.
Cauchy--Schwarz gives
\begin{equation}
\sigma_A^2\sigma_B^2\ge
|\avg{\delta A\delta B}|^2.
\end{equation}
Split the product into symmetric
and antisymmetric pieces:
\begin{equation}
\avg{\delta A\delta B}=
\frac12\avg{\{\delta A,\delta B\}}+
\frac12\avg{[A,B]}.
\end{equation}
Let $C=-\ii[A,B]$ and
$D=AB+BA-2\avg A\avg B\,I$.
Then $C,D$ are Hermitian, and
$\avg{[A,B]}=\ii\avg C$,
$\avg{\{\delta A,\delta B\}}=\avg D$.
Thus
\begin{equation}
\boxed{\sigma_A^2\sigma_B^2\ge
\frac14\left(\avg C^2+\avg D^2\right)},
\end{equation}
where $\avg C^2$ means
$(\avg C)^2$, not $\avg{C^2}$.
For $B=A$, $C=0$ and
$\avg D=2(\avg{A^2}-\avg A^2)=2\sigma_A^2$.
The extended inequality becomes
$\sigma_A^4\ge\sigma_A^4$:
it is automatically an equality and
adds no new information.

\problem{3.44}{Time evolution of a coupled three-level Hamiltonian}
In the basis $(e_1,e_2,e_3)$,
\begin{equation}
H=\begin{pmatrix}a&0&b\\0&c&0\\b&0&a\end{pmatrix}.
\end{equation}
The vector $e_2$ is an eigenstate
with energy $c$. For initial state $e_2$,
\begin{equation}
\boxed{\ket{S(t)}=\ee^{-\ii ct/\hbar}e_2}.
\end{equation}
The remaining block is
$aI+b\sigma_x$ on $\operatorname{span}(e_1,e_3)$.
Its eigenvectors $(e_1\pm e_3)/\sqrt2$
have energies $a\pm b$.
Exponentiation in that subspace gives
\begin{align}
\ee^{-\ii t(aI+b\sigma_x)/\hbar}
&=\ee^{-\ii at/\hbar}
[\cos(bt/\hbar)I-\ii\sin(bt/\hbar)\sigma_x].
\end{align}
Thus for initial state $e_1$,
\begin{equation}
\boxed{\ket{S(t)}=\ee^{-\ii at/\hbar}
\left[\cos\frac{bt}{\hbar}\,e_1
-\ii\sin\frac{bt}{\hbar}\,e_3\right]}.
\end{equation}
Its two nonzero probabilities are
$\cos^2(bt/\hbar)$ and $\sin^2(bt/\hbar)$,
summing to one.

\problem{3.45}{Position operator acting on energy coefficients}
If $\ket{S(t)}=\sum_{k=0}^\infty c_k(t)\ket k$,
insert the oscillator position matrix elements
\begin{equation}
\bra n x\ket k=
\sqrt{\frac{\hbar}{2m\omega}}
[\sqrt{k}\,\delta_{n,k-1}
+\sqrt{k+1}\,\delta_{n,k+1}].
\end{equation}
Summing over $k$ gives
\begin{equation}
\boxed{\bra n x\ket{S(t)}=
\sqrt{\frac{\hbar}{2m\omega}}
\left[\sqrt n\,c_{n-1}(t)
+\sqrt{n+1}\,c_{n+1}(t)\right]}.
\end{equation}
Set $c_{-1}=0$. The position operator
therefore mixes only neighboring energy levels;
it is not diagonal in the energy basis.

\problem{3.46}{Three-level spectra and outcome probabilities}
Let $e_1,e_2,e_3$ be the standard basis.
The specified observables are
\begin{equation}
H=\hbar\omega\,\mathrm{diag}(1,2,2),\quad
A=\lambda\begin{pmatrix}0&1&0\\1&0&0\\0&0&2\end{pmatrix},\quad
B=\mu\begin{pmatrix}2&0&0\\0&0&1\\0&1&0\end{pmatrix}.
\end{equation}
An orthonormal eigenbasis for $H$ is
$\{e_1,e_2,e_3\}$, energies
$\hbar\omega,2\hbar\omega,2\hbar\omega$.
For $A$, eigenvalues and eigenstates are
\begin{equation}
\begin{array}{c|c}
+\lambda&(e_1+e_2)/\sqrt2\\
-\lambda&(e_1-e_2)/\sqrt2\\
2\lambda&e_3
\end{array}.
\end{equation}
For $B$ they are
\begin{equation}
\begin{array}{c|c}
2\mu&e_1\\
+\mu&(e_2+e_3)/\sqrt2\\
-\mu&(e_2-e_3)/\sqrt2
\end{array}.
\end{equation}
For $\ket{S(0)}=c_1e_1+c_2e_2+c_3e_3$,
$\sum|c_j|^2=1$, the expectations are
\begin{align}
\boxed{\avg H}&=\hbar\omega
(|c_1|^2+2|c_2|^2+2|c_3|^2),\\
\boxed{\avg A}&=2\lambda
[\operatorname{Re}(c_1^*c_2)+|c_3|^2],\\
\boxed{\avg B}&=2\mu
[|c_1|^2+\operatorname{Re}(c_2^*c_3)].
\end{align}
Time evolution is
\begin{equation}
\boxed{\ket{S(t)}=
c_1\ee^{-\ii\omega t}e_1+
(c_2e_2+c_3e_3)\ee^{-2\ii\omega t}}.
\end{equation}
The energy probabilities remain
$P(\hbar\omega)=|c_1|^2$ and
$P(2\hbar\omega)=|c_2|^2+|c_3|^2$.
For $A$ the probabilities are
\begin{align}
P(A=+\lambda)&=\tfrac12|c_1\ee^{\ii\omega t}+c_2|^2,\\
P(A=-\lambda)&=\tfrac12|c_1\ee^{\ii\omega t}-c_2|^2,\\
P(A=2\lambda)&=|c_3|^2.
\end{align}
For $B$, the common phase in the
degenerate energy subspace cancels:
\begin{align}
P(B=2\mu)&=|c_1|^2,\\
P(B=+\mu)&=\tfrac12|c_2+c_3|^2,\\
P(B=-\mu)&=\tfrac12|c_2-c_3|^2.
\end{align}
Each distribution sums to one.

\problem{3.47}{Supersymmetric partner Hamiltonians and delta wells}
Write $s=\hbar/\sqrt{2m}$.
Since $p=-\ii\hbar\partial_x$,
\begin{equation}
A=s\partial_x+W(x),\quad
A^\dagger=-s\partial_x+W(x).
\end{equation}
For real $W$ and suitable domains,
multiplying and applying the product rule gives
\begin{align}
H_1=A^\dagger A&=-s^2\partial_x^2+
W^2-sW',\\
H_2=AA^\dagger&=-s^2\partial_x^2+
W^2+sW'.
\end{align}
Thus
\begin{equation}
\boxed{V_1=W^2-sW',\qquad V_2=W^2+sW'}.
\end{equation}
If $H_1\psi_n^{(1)}=E_n\psi_n^{(1)}$,
then
$H_2(A\psi_n^{(1)})=A H_1\psi_n^{(1)}
=E_n(A\psi_n^{(1)})$.
Likewise $H_1(A^\dagger\psi_n^{(2)})
=E_n(A^\dagger\psi_n^{(2)})$.
A nonzero energy eigenstate maps to a
partner state unless annihilated by the
intertwiner. If $A\psi_0^{(1)}=0$, then
\begin{equation}
\boxed{W(x)=-s\,\frac{\psi_0'(x)}{\psi_0(x)}}.
\end{equation}
For the shifted attractive delta well,
$\psi_0(x)=\sqrt\kappa\,\ee^{-\kappa|x|}$,
$\kappa=m\alpha/\hbar^2$, hence
\begin{equation}
W(x)=s\kappa\,\operatorname{sgn}(x)
=\alpha\sqrt{\frac{m}{2\hbar^2}}
\,\operatorname{sgn}(x).
\end{equation}
Use $(\operatorname{sgn}x)'=2\delta(x)$:
\begin{equation}
W^2=\frac{m\alpha^2}{2\hbar^2},\qquad
sW'=\alpha\delta(x).
\end{equation}
Therefore
\begin{equation}
\boxed{V_1=\frac{m\alpha^2}{2\hbar^2}
-\alpha\delta(x),\qquad
V_2=\frac{m\alpha^2}{2\hbar^2}
+\alpha\delta(x)}.
\end{equation}
The partner is a \emph{repulsive} delta
plus the same constant shift, so
it has no negative-energy bound level.
Its continuum scattering probabilities
equal those of the attractive partner.

\problem{3.48}{Symmetry versus self-adjointness of momentum}
Consider $p=-\ii\hbar\partial_x$ on $(0,a)$.
For differentiable $f,g$,
integration by parts gives
\begin{equation}
\braket g{pf}-\braket{pg}f
=-\ii\hbar[g^*(x)f(x)]_0^a.
\end{equation}
(a) For Dirichlet boundary conditions
$f(0)=f(a)=0$ this vanishes for every $g$
without requiring $g$ to vanish at the
endpoints. The operator is symmetric on
$H_0^1(0,a)$, but its adjoint domain is
larger (essentially $H^1(0,a)$), so it is
not self-adjoint.
(b) If $f(a)=\lambda f(0)$,
the boundary term vanishes against all
domain vectors $f$ precisely when
$\lambda g^*(a)-g^*(0)=0$,
or
\begin{equation}
\boxed{g(a)=\frac{g(0)}{\lambda^*}}
\quad(\lambda\ne0).
\end{equation}
For the adjoint domain to be identical
to the operator domain, require
$\lambda=1/\lambda^*$, equivalently
$\boxed{|\lambda|=1}$.
Thus all phases $\lambda=\ee^{\ii\theta}$
give self-adjoint quasi-periodic momentum
operators on the finite interval.
(c) On $(0,\infty)$, the deficiency equations
$p^\dagger g=\pm\ii\gamma g$, $\gamma>0$,
give
$g_+(x)\propto\ee^{-\gamma x/\hbar}$,
$g_-(x)\propto\ee^{+\gamma x/\hbar}$.
Only the first is square integrable,
so the deficiency indices are unequal:
$\boxed{(n_+,n_-)=(1,0)}$ for this convention.
Therefore there is no self-adjoint
extension of $p$ on the half-line.
\qnote{Momentum domains in a finite box
are not determined solely by the
formal differential expression.}

\problem{3.49}{Gaussian free-particle evolution in momentum space}
(a) For $H=p^2/(2m)$ the momentum-space
Schrödinger equation is an ODE at each $p$:
\begin{equation}
\ii\hbar\partial_t\Phi(p,t)=
\frac{p^2}{2m}\Phi(p,t)
\quad\Longrightarrow\quad
\boxed{\Phi(p,t)=
\Phi(p,0)\ee^{-\ii p^2t/(2m\hbar)}}.
\end{equation}
(b) For the traveling Gaussian
$\Psi(x,0)=(2a/\pi)^{1/4}
\ee^{-ax^2+\ii\ell x}$ with $a>0$,
completing the Fourier-transform square yields
\begin{equation}
\boxed{\Phi(p,0)=
(2\pi a\hbar^2)^{-1/4}
\exp\left[-\frac{(p-\hbar\ell)^2}
{4a\hbar^2}\right]}.
\end{equation}
Therefore $|\Phi(p,t)|^2$ is the
time-independent Gaussian
\begin{equation}
\boxed{|\Phi(p,t)|^2=
\frac1{\sqrt{2\pi a\hbar^2}}
\exp\left[-\frac{(p-\hbar\ell)^2}
{2a\hbar^2}\right]}.
\end{equation}
(c) Standard Gaussian integration gives
\begin{equation}
\boxed{\avg p=\hbar\ell,\quad
\avg{p^2}=\hbar^2(\ell^2+a),\quad
\sigma_p=\hbar\sqrt a}.
\end{equation}
(d) The energy expectation decomposes as
\begin{equation}
\boxed{\avg H=\frac{\hbar^2(\ell^2+a)}{2m}
=\frac{\avg p^2}{2m}
+\underbrace{\frac{\hbar^2 a}{2m}}_{\avg H_0}},
\end{equation}
where $\avg p^2$ in the first fraction
of the final expression denotes
$(\avg p)^2$, not $\avg{p^2}$.
The first term is center-of-mass
translation energy and the second is
the internal spreading energy of
the Gaussian wavepacket.

\section{Reproducible figures}
The companion script \texttt{scripts/plots.py} produces
momentum-density and coherent-state illustrations.
\IfFileExists{figures/momentum_well.pdf}{
\begin{figure}[htbp]\centering
\includegraphics[width=.86\linewidth]{figures/momentum_well.pdf}
\caption{Momentum distributions for infinite-well states $n=1,2,5,10$
(Problem 3.34), numerically normalized and compared with the exact transform.}
\end{figure}}{}
\IfFileExists{figures/coherent.pdf}{
\begin{figure}[htbp]\centering
\includegraphics[width=.86\linewidth]{figures/coherent.pdf}
\caption{Coherent harmonic-oscillator probability density translated in
position without changing its width (Problem 3.42).}
\end{figure}}{}
\clearpage
\begin{thebibliography}{9}
\bibitem{griffiths} D.~J.~Griffiths and D.~F.~Schroeter,
\emph{Introduction to Quantum Mechanics}, 3rd ed.,
Cambridge University Press (2018), Chapter 3, pp. 91–130.
DOI: \href{https://doi.org/10.1017/9781316995433}{10.1017/9781316995433}.
\end{thebibliography}
\end{document}